\chapter{Chapter 1: A Framework for LLM System Design Interviews}

\section*{Context}
\textbf{Problem statement:} Build a repeatable framework for answering any LLM system design interview question with production-grade technical depth in 40-45 minutes.


\subsection*{Real-World Examples}
\begin{itemize}
  \item A fintech team designs an internal LLM assistant and must justify why RAG is safer than fine-tuning for policy updates every week.
  \item A healthcare startup interview asks for a triage bot design where abstention and escalation logic matter more than fluent prose.
\end{itemize}

\section*{Step 1: Understand the Problem and Establish Design Scope}
\subsection*{Simulated Dialogue}
\textbf{Interviewer:} Design a framework for LLM system design interviews.\\
\textbf{Candidate:} I will scope product goals, risk profile, and numeric constraints first, then map them to architecture and trade-offs.

\subsection*{Functional Requirements}
\begin{itemize}
  \item Provide a reusable 5-step response structure for any LLM system design prompt.
  \item Force explicit design choices for retrieval, model routing, and safety controls.
  \item Include deterministic fallback paths for low-confidence and high-risk states.
  \item Produce outputs that are implementable and measurable.
\end{itemize}

\subsection*{Non-Functional Requirements}
\begin{itemize}
  \item Interview execution: coherent high-level design in under 10 minutes.
  \item Latency discipline with explicit TTFT and TPOT budgets.
  \item Cost visibility via token economics and serving envelope.
  \item Scalability, safety, and rollback readiness.
\end{itemize}

\subsection*{Back-of-the-envelope Estimation}
\begin{itemize}
  \item DAU: 1.5M, Requests/day: 18M, Avg input: 1,600 tokens, Avg output: 280 tokens.
  \item Average QPS: $18,000,000/86,400 \approx 208$, Peak QPS: $\approx 2,500$.
  \item Input tokens/day: $28.8B$, Output tokens/day: $5.04B$.
  \item Daily API cost: $(input/10^6)\cdot p_{in} + (output/10^6)\cdot p_{out}$.
  \item Vector storage example (300M chunks, 1536-dim float32): raw $\approx 1.8TB$, production $\approx 4.0$--$5.5TB$ with overhead and replicas.
  \item VRAM planning (order-of-magnitude): 7B quantized $\sim 16$--$24GB$, 13B--34B $\sim 40$--$80GB$, 70B class requires multi-GPU parallel serving.
\end{itemize}

\section*{Step 2: Propose High-Level Design and Get Buy-In}
\subsection*{Core Architecture Components}
\begin{itemize}
  \item Client + gateway, policy layer, orchestrator, retrieval stack, model router.
  \item Serving runtime with continuous batching, KV/prefix cache, and streaming.
  \item Cache hierarchy: retrieval cache, semantic cache, response cache.
  \item Safety stack and observability/evaluation loop.
\end{itemize}

\subsection*{Pipeline A: Data Ingestion / Indexing}
\begin{enumerate}
  \item Ingest and normalize sources.
  \item Parse and chunk (including parent-child when needed).
  \item Embed and index vectors + lexical terms + metadata.
  \item Validate index quality and promote by alias swap.
\end{enumerate}

\subsection*{Pipeline B: Inference / Query}
\begin{enumerate}
  \item Authenticate and apply quota/policy checks.
  \item Classify intent and choose execution path (direct/RAG/tool/agent).
  \item Retrieve, rerank, pack context, and route model tier.
  \item Stream output, run safety/grounding verification, emit traces.
\end{enumerate}

\textit{[Insert Architecture Diagram: Show gateway -> policy layer -> orchestrator branching into direct LLM, RAG, and tool/agent paths; include hybrid retrieval + reranker, model router with managed/self-hosted tiers, continuous batching runtime with KV cache, semantic cache check, and observability feedback loop feeding evaluation and release gates.]}

\subsection*{TikZ Architecture Diagram}
\begin{center}
\begin{tikzpicture}[node distance=1.6cm, auto, >=stealth]
  \node[draw, rounded corners] (client) {UserClient};
  \node[draw, rounded corners, right=of client] (gateway) {APIGateway};
  \node[draw, rounded corners, right=of gateway] (orch) {Orchestrator};
  \node[draw, rounded corners, above right=0.8cm and 1.6cm of orch] (retr) {HybridRetrieval};
  \node[draw, rounded corners, below right=0.8cm and 1.6cm of orch] (mr) {ModelRouter};
  \node[draw, rounded corners, right=3.2cm of orch] (run) {ServeRuntime};
  \node[draw, rounded corners, below=of run] (safe) {SafetyVerifier};
  \draw[->] (client) -- (gateway);
  \draw[->] (gateway) -- (orch);
  \draw[->] (orch) -- (retr);
  \draw[->] (orch) -- (mr);
  \draw[->] (retr) -- (mr);
  \draw[->] (mr) -- (run);
  \draw[->] (run) -- (safe);
  \draw[->] (safe) |- (client);
\end{tikzpicture}
\end{center}

\section*{Step 3: Design Deep Dive}
\begin{enumerate}
  \item \textbf{Interview Scoping for LLM Systems}
  \item \textbf{Cost-Latency-Quality-Safety Trade-off Engine}
  \item \textbf{Production Failure Modes and Mitigations}
\end{enumerate}

\subsection*{Trade-off Comparisons}
\begin{tabular}{|p{0.2\linewidth}|p{0.24\linewidth}|p{0.24\linewidth}|p{0.24\linewidth}|}
\hline
\textbf{Component} & \textbf{Option A} & \textbf{Option B} & \textbf{Recommended} \\
\hline
Retrieval & Keyword-only & Dense-only & Hybrid + cross-encoder reranking \\
\hline
Chunking & Fixed-size & Semantic & Parent-child for large docs \\
\hline
Serving & Managed-only & Self-hosted-only & Hybrid managed + self-hosted \\
\hline
Caching & No cache & Semantic cache only & Semantic + KV cache hierarchy \\
\hline
\end{tabular}

\subsection*{Production Bottlenecks}
\begin{itemize}
  \item Provider rate limits and retry storms.
  \item Queue saturation and head-of-line blocking.
  \item GPU OOM from long prompts and oversized dynamic batches.
  \item Stale embeddings and index drift.
  \item Metadata/ACL connection pool exhaustion.
  \item Agentic hallucination loops without step budgets.
\end{itemize}

\subsection*{Advanced Techniques to Mention Explicitly}
\begin{itemize}
  \item KV caching
  \item Continuous batching
  \item Semantic caching
  \item Cross-encoder reranking
  \item Parent-child chunking
  \item LoRA adapters
\end{itemize}

\section*{Step 4: Wrap Up}
\textbf{Final design:} A deterministic decision framework that prioritizes quantification first, then routes execution by complexity/risk, and enforces safety plus rollback as first-class concerns.

\textbf{Wrap-around concerns:}
\begin{itemize}
  \item Release engineering with canary/shadow/rollback triggers.
  \item Governance with policy versioning and audit-complete traces.
  \item Route-level cost controls and drift monitoring by version.
\end{itemize}

\section*{Step 5: Interviewer Curveballs}
\subsection*{Interviewer Follow-ups \& Frequently Asked Questions}
\textbf{Q1:} TTFT regressed from 900ms to 2.4s overnight. Where first?\\
\textbf{The Ideal Answer:} Check queue wait, cache-hit regressions, then provider throttling and scheduler saturation.\\
\textbf{Deep Dive \& Trade-offs:} Prioritize queue metrics and route distribution; apply emergency token caps and fast-tier rerouting.\\
\textbf{The Interviewer's Rubric:} Tests incident triage under LLM latency pressure.

\textbf{Q2:} Managed vs self-hosted for one route?\\
\textbf{The Ideal Answer:} Choose by all-in unit economics and SLO risk at steady-state volume.\\
\textbf{Deep Dive \& Trade-offs:} Managed for speed, self-hosted for stable high-throughput economics, hybrid for resilience.\\
\textbf{The Interviewer's Rubric:} Tests pragmatic hosting decisions.

\textbf{Q3:} Default token budget policy?\\
\textbf{The Ideal Answer:} Hard route-level input/output caps with context-packing quotas and truncation policy.\\
\textbf{Deep Dive \& Trade-offs:} Prevents OOM, tail latency, and cost spikes while preserving relevance via reranking.\\
\textbf{The Interviewer's Rubric:} Tests context-window operationalization.

\textbf{Q4:} Avoid stale vector quality without full re-embedding?\\
\textbf{The Ideal Answer:} Incremental re-embedding plus periodic shadow-index evaluation and alias-based promotion.\\
\textbf{Deep Dive \& Trade-offs:} Balances freshness and cost; keep rollback alias hot.\\
\textbf{The Interviewer's Rubric:} Tests lifecycle rigor.

\textbf{Q5:} What breaks first at 10x spikes?\\
\textbf{The Ideal Answer:} Queueing and upstream rate limits, then cache misses amplify model load.\\
\textbf{Deep Dive \& Trade-offs:} Use reserved capacity lanes, deterministic degrade modes, and bounded retries.\\
\textbf{The Interviewer's Rubric:} Tests burst resilience design.

\textbf{Q6:} Stop agent tool loops?\\
\textbf{The Ideal Answer:} Max-step budgets, loop-signature detection, and forced terminal action.\\
\textbf{Deep Dive \& Trade-offs:} Enforce task-class step ceilings and escalation fallback.\\
\textbf{The Interviewer's Rubric:} Tests control-plane safety.

\textbf{Q7:} Minimum weekly health metrics?\\
\textbf{The Ideal Answer:} TTFT p95, TPOT p95, cost/query, Recall@K, citation precision, policy violation rate.\\
\textbf{Deep Dive \& Trade-offs:} Monitor quality, cost, and safety jointly by route and version.\\
\textbf{The Interviewer's Rubric:} Tests multidimensional operating discipline.

\textbf{Q8:} Why mention LoRA adapters in interviews?\\
\textbf{The Ideal Answer:} They deliver low-cost domain adaptation with faster iteration than full fine-tuning.\\
\textbf{Deep Dive \& Trade-offs:} Adds routing/version complexity; keep base-model rollback path.\\
\textbf{The Interviewer's Rubric:} Tests practical adaptation knowledge.

\section*{End of Chapter}
\begin{itemize}
  \item \textbf{Key takeaway:} Tie architecture decisions to measurable SLOs, cost, and safety outcomes.
  \item \textbf{What to practice:} Defend one default and one fallback design with explicit trade-offs.
  \item \textbf{Interview drill:} Explain regression detection, rollback triggers, and risk communication.
\end{itemize}
